% from ega4-18-fr.tex
\oldpage[IV-4]{118}

\begin{lemma}[18.4.1]
\label{IV.18.4.1}
Soient $A$ un anneau, $B$ une $A$-algèbre finie monogène, $u$ un générateur de la $A$-algèbre $B$, $F\in A[T]$ un polynôme tel que $F(u)=0$, $F'$ le polynôme dérivé; posons $u'=F'(u)$. Alors l'idéal de $B$, annulateur de $\Omega^1_{B/A}$, contient $u'B$; il est égal à $u'B$ si l'idéal $\mathfrak J$ de $A[T]$ formé des polynômes $G$ tels que $G(u)=0$, est engendré par $F$, autrement dit si l'homomorphisme surjectif canonique $\varphi:A[T]/F.A[T]\to B$ transformant l'image de $T$ en $u$, est bijectif.
\end{lemma}

Posons $C=A[T]$, de sorte que $B=C/\mathfrak J$. On a la suite exacte (0, 20.5.12.1)
\[
  \mathfrak J/\mathfrak J^2\longrightarrow\Omega^1_{C/A}\otimes_C B\longrightarrow\Omega^1_{B/A}\longrightarrow 0
\]
et $\Omega^1_{B/A}$ s'identifie donc au quotient $B/\mathfrak J'$, $\mathfrak J'$ étant l'idéal engendré par les éléments $G'(u)$, où $G$ parcourt un système de générateurs de l'idéal $\mathfrak J$ (0, 20.5.13); d'où aussitôt le lemme.

\begin{proposition}[18.4.2]
\label{IV.18.4.2}
Les notations étant celles de (18.4.1), soit $\mathfrak q$ un idéal premier de $B$. Alors:
\begin{enumerate}
  \item[(i)] Si $\mathfrak q$ ne contient pas $u'$, $B_{\mathfrak q}$ est une $A_{\mathfrak p}$-algèbre formellement non ramifiée ($\mathfrak p$ étant l'image réciproque de $\mathfrak q$ dans $A$); en d'autres termes, $\Spec(B_{u'})$ est formellement non ramifié sur $\Spec(A)$.
  \item[(ii)] Supposons de plus que $F$ soit unitaire et engendre $\mathfrak J$. Alors, pour que $\Spec(B)$ soit étale sur $\Spec(A)$ au point $\mathfrak q$, il faut et il suffit que $u'\notin\mathfrak q$.
\end{enumerate}
\end{proposition}

L'hypothèse que $u'\notin\mathfrak q$ entraîne que $\Omega^1_{B_{\mathfrak q}/A_{\mathfrak p}}=0$ (0, 20.5.9), donc (i) résulte de (17.2.1). En outre, sous les hypothèses de (ii), $B$ est un $A$-module libre en vertu de la division euclidienne; comme l'annulateur $\mathfrak J'$ de $\Omega^1_{B/A}$ est alors égal à $u'B$ en vertu de (18.4.1) et que $\Omega^1_{B/A}$ est un $B$-module de présentation finie (16.4.22), l'annulateur de $\Omega^1_{B_{\mathfrak q}/A_{\mathfrak p}}$ est égal à $u'B_{\mathfrak q}$ (Bourbaki, \emph{Alg. comm.}, chap. II, \textsection~2, n\textsuperscript{o}~4, formule (9)), et (ii) résulte donc de (i) et de l'implication $c)\Rightarrow a)$ dans (17.6.1).

\begin{corollary}[18.4.3]
\EGAReaderTarget{ega4.c04.cor.18-4-3}{IV.18.4.3}
Les notations étant celles de (18.4.2), supposons que $F$ soit unitaire et engendre $\mathfrak J$. Alors, pour que $B$ soit une $A$-algèbre étale, il faut et il suffit que $u'$ soit inversible dans $B$ (ou, ce qui revient au même, que l'idéal de $A[T]$ engendré par $F$ et $F'$ soit égal à $A[T]$).
\end{corollary}

Compte tenu de (18.4.2), (ii), dire que $\Spec(B)$ est étale sur $\Spec(A)$ signifie en effet que $u'$ n'appartient à aucun idéal premier de $B$, i.e. qu'il est inversible dans $B$.

On dit qu'un polynôme unitaire $F\in A[T]$ tel que l'idéal de $A[T]$ engendré par $F$ et $F'$ soit égal à $A[T]$ lui-même, est \emph{séparable}; il est immédiat que cette définition coïncide avec la définition usuelle (Bourbaki, \emph{Alg.}, chap. V, \textsection~7, n\textsuperscript{o}~6) lorsque $A$ est un corps.

\begin{lemma}[18.4.4]
\label{IV.18.4.4}
Soient $A$ un anneau local, $B$ une $A$-algèbre finie monogène, $u$ un générateur de la $A$-algèbre $B$. Soient $\mathfrak n_i$ ($1\leq i\leq r$) des idéaux maximaux de $B$ tels que $\Spec(B)$ soit formellement non ramifié sur $\Spec(A)$ aux points $\mathfrak n_i$. Alors il existe un polynôme unitaire $F\in A[T]$ tel que $F(u)=0$ et $F'(u)\notin\mathfrak n_i$ pour tout indice $i$ ($1\leq i\leq r$). En outre, si $k$ est le corps

